
\documentclass{article}
\usepackage{colortbl}
\usepackage{makecell}
\usepackage{multirow}
\usepackage{supertabular}

\begin{document}

\newcounter{utterance}

\centering \large Interaction Transcript for game `cladder', experiment `full\_v1.5\_default', episode 9213 with Qwen3.6{-}27B.
\vspace{24pt}

{ \footnotesize  \setcounter{utterance}{1}
\setlength{\tabcolsep}{0pt}
\begin{supertabular}{c@{$\;$}|p{.15\linewidth}@{}p{.15\linewidth}p{.15\linewidth}p{.15\linewidth}p{.15\linewidth}p{.15\linewidth}}
    \# & \multicolumn{2}{c}{Player} && \multicolumn{2}{c}{Game Master} \\
    \hline

    \theutterance \stepcounter{utterance}  
    & & & \multicolumn{4}{p{0.6\linewidth}}{
        \cellcolor[rgb]{0.9,0.9,0.9}{
            \makecell[{{p{\linewidth}}}]{
                \texttt{\tiny{[P1$\langle$GM]}}
                \texttt{You are an expert in causal inference. The following question is not a typical commonsense query, but rather a meticulously designed question created by a professor specializing in causal inference, intended to assess the students' mastery of the course content.} \\
\texttt{Imagine a self{-}contained, hypothetical world with only the following conditions, and without any unmentioned factors or causal relationships: Pexu has a direct effect on rukz. Hwax has a direct effect on rukz. For those who are not pexu, the probability of rukz is 9\%. For those who are pexu, the probability of rukz is 18\%. For those who are pexu, would it be less likely to see rukz if the individual was not pexu?} \\
\texttt{Start your answer with 'Yes' or 'No', followed by additional reasoning or evidence to support your explanation.} \\
            }
        }
    }
     \\ \\

    \theutterance \stepcounter{utterance}  
    & \multicolumn{4}{p{0.6\linewidth}}{
        \cellcolor[rgb]{0.561,0.737,0.561}{
            \makecell[{{p{\linewidth}}}]{
                \texttt{\tiny{[P1$\rangle$GM]}}
                \texttt{Yes} \\
\\ 
\texttt{**Reasoning:**} \\
\\ 
\texttt{To answer this question, we must interpret the causal structure and the specific counterfactual query provided.} \\
\\ 
\texttt{1.  **Identify the Causal Structure:**} \\
\texttt{    *   \$Pexu \textbackslash{}rightarrow Rukz\$ (Direct effect)} \\
\texttt{    *   \$Hwax \textbackslash{}rightarrow Rukz\$ (Direct effect)} \\
\texttt{    *   There are no other factors or relationships. This implies that \$Pexu\$ and \$Hwax\$ are independent of each other (no confounding between them mentioned, and no unmentioned factors).} \\
\\ 
\texttt{2.  **Analyze the Probabilities:**} \\
\texttt{    *   \$P(Rukz | \textbackslash{}neg Pexu) = 0.09\$ (9\%)} \\
\texttt{    *   \$P(Rukz | Pexu) = 0.18\$ (18\%)} \\
\\ 
\texttt{3.  **Interpret the Question:**} \\
\texttt{    The question asks: *"For those who are pexu, would it be less likely to see rukz if the individual was not pexu?"*} \\
\\ 
\texttt{    This is a counterfactual question regarding the **Individual{-}Level Causal Effect** or the **Average Causal Effect** for the treated group. Specifically, it asks us to compare the observed outcome for an individual who is currently \$Pexu\$ (\$Rukz\$) with the potential outcome for that same individual if they had been \$\textbackslash{}neg Pex} \\
            }
        }
    }
    & & \\ \\

    \theutterance \stepcounter{utterance}  
    & & & \multicolumn{2}{p{0.3\linewidth}}{
        \cellcolor[rgb]{0.9,0.9,0.9}{
            \makecell[{{p{\linewidth}}}]{
                \texttt{\tiny{[GM$|$GM]}}
                \texttt{Yes} \\
            }
        }
    }
    & & \\ \\

    \theutterance \stepcounter{utterance}  
    & & & \multicolumn{2}{p{0.3\linewidth}}{
        \cellcolor[rgb]{0.9,0.9,0.9}{
            \makecell[{{p{\linewidth}}}]{
                \texttt{\tiny{[GM$|$GM]}}
                \texttt{yes} \\
            }
        }
    }
    & & \\ \\

    \theutterance \stepcounter{utterance}  
    & & & \multicolumn{2}{p{0.3\linewidth}}{
        \cellcolor[rgb]{0.9,0.9,0.9}{
            \makecell[{{p{\linewidth}}}]{
                \texttt{\tiny{[GM$|$GM]}}
                \texttt{game\_result = WIN} \\
            }
        }
    }
    & & \\ \\

\end{supertabular}
}

\end{document}
